\subsection*{真实出现频率与编码假设可以不同}
假设真实来源是公平二元分布，但编码或预测模型认为第一类占四分之三。模型因此给第一类较短的理想长度，给第二类较长的长度；实际使用频率仍由真实分布决定。平均时用真实概率加权模型码长，得到交叉熵。

这解释了公式里为什么同时出现 $p$ 和 $q$：$p$ 决定现实中哪种事件出现多少次，$q$ 决定每次出现时付多大代价。若把两个位置交换，就换了问题，不能期待相同答案。

\subsection*{把总代价分成不可避免与模型失配}
在每项里加减 $\ln p_i$，交叉熵分解为真实熵加 KL 散度。前者是来源自身的不确定性，后者衡量采用模型分布多出的平均对数代价。在两分布归一化且严格为正时，用 $\ln u\le u-1$ 可证明额外代价非负。

在交互实验中固定真实概率 $p$，只移动模型概率 $q$。预测交叉熵在哪里最低，再观察 $q=p$ 时曲线达到真实熵。把 $q$ 推近零或一，会看到模型对实际可能事件过度否定时，代价迅速增加。

分类训练中，单样本真实标签只在正确类别上取一，所以损失变成 $-\ln q_j$。这让“提高正确类概率”与“减少编码代价”连接起来。但训练数据上的低损失不自动保证新数据上的校准或可靠性。数学分解、统计估计与模型泛化仍需分别论证。
